% SPDX-License-Identifier: CC-BY-NC-ND-4.0
% Copyright 2026 Ingolf Lohmann.
\documentclass[11pt,a4paper]{article}

\usepackage{fontspec}
\setmainfont{Latin Modern Roman}
\setsansfont{Latin Modern Sans}
\setmonofont{Latin Modern Mono}[Scale=0.88]
\usepackage[provide=*,main=ngerman,english]{babel}
\usepackage{microtype}
\usepackage{geometry}
\usepackage{amsmath,amssymb,mathtools,amsthm}
\usepackage{booktabs,longtable,tabularx,array}
\usepackage{enumitem}
\usepackage{xcolor}
\usepackage{fancyhdr}
\usepackage{fancyvrb}
\usepackage{seqsplit}
\usepackage{xurl}
\usepackage{hyperref}
\usepackage{bookmark}
\usepackage[nameinlink,noabbrev]{cleveref}
\newcolumntype{L}[1]{>{\raggedright\arraybackslash}p{#1}}

\geometry{left=24mm,right=24mm,top=24mm,bottom=27mm,headheight=14pt}
\setlength{\parindent}{0pt}
\setlength{\parskip}{0.55em}
\setlength{\emergencystretch}{3em}
\setlist[itemize]{leftmargin=1.45em,itemsep=0.18em,topsep=0.3em}
\setlist[enumerate]{leftmargin=1.7em,itemsep=0.22em,topsep=0.3em}

\definecolor{QBlue}{HTML}{153E69}
\definecolor{QLight}{HTML}{ECF3F8}
\definecolor{QGreen}{HTML}{E8F3EA}
\definecolor{QAmber}{HTML}{FFF2CC}
\definecolor{QRed}{HTML}{FBE9E7}
\definecolor{QGray}{HTML}{4D5660}

\hypersetup{
  colorlinks=true,
  linkcolor=QBlue,
  citecolor=QBlue,
  urlcolor=QBlue,
  pdftitle={QIK-VRT: Wirkungsgeprüfte bidirektionale Kommunikation über virtuelle Zeit},
  pdfauthor={Ingolf Lohmann},
  pdfsubject={Ausführbarer C90-Zeuge, bedingte Sätze und die Grenze zur physikalischen Retrokausalität},
  pdfkeywords={QIK-VRT, virtuelle Zeit, bidirektionale Kommunikation, Replay, Effect Acknowledgement, Informationsübertragung, Retrokausalität}
}

\pagestyle{fancy}
\fancyhf{}
\fancyhead[L]{\small\textcolor{QGray}{QIK--VRT: virtueller Zeitkanal}}
\fancyhead[R]{\small\textcolor{QGray}{Ingolf Lohmann}}
\fancyfoot[C]{\small\thepage}
\renewcommand{\headrulewidth}{0.3pt}

\newtheorem{definition}{Definition}[section]
\newtheorem{satz}{Satz}[section]
\newtheorem{korollar}{Korollar}[section]
\newtheorem{lemma}{Lemma}[section]
\theoremstyle{remark}
\newtheorem{bemerkung}{Bemerkung}[section]
\crefname{definition}{Definition}{Definitionen}
\crefname{satz}{Satz}{Sätze}
\crefname{korollar}{Korollar}{Korollare}
\crefname{lemma}{Lemma}{Lemmata}

\newcommand{\qikvrt}{QIK--VRT}
\newcommand{\DONE}{\texttt{EFFECT\_ACK\_DONE}}
\newcommand{\bytes}{\{0,1\}^{*}}
\newcommand{\hash}[1]{{\ttfamily\small\seqsplit{#1}}}
\newcommand{\statusbox}[2]{%
  \par\medskip\noindent
  \colorbox{#1}{\parbox{\dimexpr\linewidth-2\fboxsep\relax}{#2}}%
  \par\medskip
}

\begin{document}

\hypersetup{pageanchor=false}
\begin{titlepage}
  \thispagestyle{empty}
  \vspace*{12mm}
  {\color{QBlue}\rule{\textwidth}{1.4pt}}\\[1.0em]
  {\Huge\bfseries QIK--VRT:\\[0.18em]
  Der bidirektionale virtuelle Zeitkanal\par}
  \vspace{0.9em}
  {\LARGE Ausführbarer C90-Zeuge, bedingte Sätze und\\
  die Grenze zur physikalischen Retrokausalität\par}
  \vspace{1.0em}
  {\color{QBlue}\rule{\textwidth}{0.6pt}}

  \vfill
  {\large\textbf{Ingolf Lohmann}\par}
  \vspace{0.35em}
  {\normalsize Independent Researcher; QIK--VRT\par}
  \vspace{0.35em}
  {\normalsize Working Paper, Version 1.0.0 -- 1. August 2026\par}

  \vfill
  \statusbox{QLight}{
    \textbf{Tragender Satz.}
    Ein ausführbares System kann vollständige endliche Dialoge über
    rückwärtsgerichtete virtuelle Adressen führen, ohne seine physische
    Hostkausalität oder seine unveränderliche Quellhistorie zu verletzen.
    Der neue ISO-C90-Zeuge führt Anfrage und Antwort bytegenau in beiden
    virtuellen Richtungen aus. Der allgemeine Satz für beliebige endliche
    Bitfolgen wird konstruktiv bewiesen; sein frischer Lean-Kerneltransfer ist
    noch offen. Eine physikalische Zukunft-zu-Vergangenheit-Kanalkapazität
    folgt daraus nicht und wird als getrennte experimentelle Hypothese
    behandelt.
  }

  \vfill
  {\small
    Dokumentationslizenz: CC BY-NC-ND 4.0.\\
    Vorveröffentlichungskandidat; nicht peer-reviewed; kein neuer DOI.\\
    Das zitierte Internet-Draft ist ein individuelles Arbeitsdokument und
    keine IETF-Billigung.
  }
\end{titlepage}
\hypersetup{pageanchor=true}
\setcounter{page}{1}

\selectlanguage{ngerman}

\begin{abstract}
Diese Arbeit trennt drei Ordnungen, die in Debatten über digitale
Retrokausalität häufig vermischt werden: die strikt vorwärts laufende
Hostordnung, die frei adressierbare virtuelle Zeitordnung und die gesondert
freizugebende Effektordnung. Auf dieser Grundlage wird ein
bidirektionaler virtueller Zeitkanal definiert. Eine später erzeugte endliche
Nachricht adressiert einen früheren, unveränderlich identifizierten
virtuellen Zustand; ein deterministischer Branch-Replay verarbeitet sie und
sendet eine gebundene Antwort an eine gegenwärtige virtuelle Adresse zurück.

Ein neuer, dependency-freier ISO-C90-Zeuge realisiert diesen Ablauf mit neun
streng aufsteigenden Hostereignissen. Die Detailausführung überträgt 257 Byte
in 16 Blöcken von der virtuellen Adresse 30 zur Adresse 15 und 258 Byte zurück.
Zwei unabhängige Replays erzeugen dieselbe Antwort, der Quell-Digest bleibt
unverändert, zehn Grenzlängen von 0 bis 4096 Byte bestehen in beiden
Richtungen, und ein absichtlich fehlender Block wird verworfen. Diese
Ausführung ist reproduzierbare Evidenz für einen endlichen Softwarezeugen,
kein universeller Programmbeweis.

Konstruktiv wird gezeigt: Für jede endliche Bitfolge existiert unter
integrer, schließlich erfolgreicher Zustellung und ausreichenden endlichen
Ressourcen eine bytegenaue Segmentierung und Reassemblierung. Zwei solcher
Kanäle, eine terminierende Antwortfunktion, eindeutige Sessionbindungen und
vollständig mediierte Wirkungsgates ergeben einen endlichen bidirektionalen
Dialog. Alle erzeugenden Hostereignisse bleiben dabei vorwärtsgerichtet; die
Basishistorie wird nicht überschrieben. Der bestehende Lean-4-Kern von
\qikvrt{} bleibt als maschinengeprüfte Grundlage für DONE-only-Freigabe,
kontrafaktische Zukunftsrelevanz im endlichen Modell, Nichtüberschreibung der
Vergangenheitsprojektion und reziproke Receipt-Bindung erhalten. Die neuen
Kompositionssätze sind papierbewiesen, aber noch nicht durch einen frischen
Kernel-Receipt klassifiziert.

Die Einzelmechanismen besitzen erhebliche Vorarbeiten in logischer und
virtueller Zeit, bitemporalen Datenbanken, retroaktiven Datenstrukturen,
Event Sourcing, End-to-End-Argumenten, Effektsystemen und Proof-Carrying
Authorization. Der scope-gebundene Neuheitskandidat ist ihre explizite
Verbindung mit unveränderlicher Quellhistorie, Branch-Replay, drei getrennten
Ordnungen, bidirektionaler Sessionkorrelation und fail-closed
Wirkungsautorisierung. Für eine Behauptung physikalischer
Rückwärtssignalisierung wäre zusätzlich eine operative Brücke von virtuellen
Adressen zu Raumzeitereignissen sowie positive, reproduzierbare
Kanalkapazität von einer später randomisierten Intervention zu einem zuvor
versiegelten Messwert erforderlich. Dafür liefert der Softwarezeuge keine
Evidenz.
\end{abstract}

\begin{otherlanguage}{english}
\begin{abstract}
This paper separates three orders often conflated in discussions of digital
retrocausality: strictly forward host order, freely addressable virtual time,
and separately authorized effect order.  It defines a bidirectional virtual
time channel in which a later-created finite message targets an immutable
earlier virtual state; deterministic branch replay processes the message and
returns a bound response to a present virtual address.

A new dependency-free ISO C90 witness executes the construction using nine
strictly increasing host events.  Its detailed run transports 257 bytes in
16 chunks from virtual address 30 to address 15 and 258 bytes back.  Two
independent replays return the same answer, the source digest remains
unchanged, ten boundary lengths from 0 through 4096 bytes pass in both
directions, and a deliberately omitted chunk is rejected.  This is
reproducible evidence for one finite executable model, not a universal
program proof.

We prove constructively that every finite bit string can be segmented and
reassembled byte-exactly under integrity, eventual delivery, and sufficient
finite resources.  Two such channels, a terminating response function,
unique session bindings, and completely mediated effect gates yield a finite
bidirectional dialogue while every generating host event remains forward
ordered and the base history remains unchanged.  Existing Lean 4 receipts
continue to establish the finite QIK--VRT release and reciprocal-binding
kernel.  The new composition theorems have complete paper proofs, but no new
kernel receipt yet.  No physical backwards-signalling claim follows.  Such a
claim would require an additional operational bridge to spacetime events and
positive reproducible channel capacity from a later randomized intervention
to an already sealed earlier record.
\end{abstract}
\end{otherlanguage}

\setcounter{tocdepth}{1}
\tableofcontents
\newpage

\section{Ergebnis und Anspruchsordnung}

Der zentrale Informatiksatz dieses Dokuments lautet nicht, dass ein realer
Prozessor rückwärts durch die Raumzeit läuft. Er lautet:

\statusbox{QGreen}{
\textbf{Virtueller Dialogsatz in Kurzform.}
Ein gegenwärtiger Hostprozess kann eine endliche Nachricht an einen
unveränderlich bezeichneten früheren virtuellen Zustand adressieren, dort
einen deterministischen Branch berechnen und dessen endliche Antwort
bytegenau an eine gegenwärtige virtuelle Adresse zurückführen. Anfrage und
Antwort bilden einen bidirektionalen virtuellen Informationsdialog, während
sämtliche tatsächlichen Rechen- und Wirkungsereignisse streng vorwärts
geordnet bleiben.
}

Damit ein stark formulierter Satz nicht unbemerkt seine Beweisart wechselt,
verwendet das Dokument fünf epistemische Klassen:

\begin{longtable}{@{}L{0.20\linewidth}L{0.72\linewidth}@{}}
\toprule
\textbf{Klasse} & \textbf{Bedeutung in dieser Arbeit} \\
\midrule
\endhead
\texttt{EXECUTED} & Ein konkretes Programm wurde in der benannten Toolchain
gebaut und ausgeführt; der Befund gilt für diesen endlichen Lauf.\\
\texttt{KERNEL\_VERIFIED} & Ein benannter formaler Satz wurde im bestehenden
Lean-4-Scope mit leerer Axiomenliste geprüft und durch einen Receipt gebunden.\\
\texttt{CONDITIONAL} & Ein mathematischer Satz wird hier aus expliziten
Voraussetzungen bewiesen; ein neuer Kernel-Receipt steht noch aus.\\
\texttt{ENGINEERING} & Eine konstruktive Protokollerweiterung ist spezifiziert,
aber noch nicht unabhängig interoperabel implementiert.\\
\texttt{OPEN\_PHYSICAL} & Eine physikalische Behauptung benötigt eine
zusätzliche Brücke und neue Messdaten; sie folgt nicht aus dem Softwaremodell.\\
\bottomrule
\end{longtable}

Diese Ordnung ist keine rhetorische Abschwächung. Sie verhindert den
Fehlschluss, aus einem richtigen Informatiksatz eine andere physikalische
Behauptung zu gewinnen. Zugleich verhindert sie die umgekehrte
Untertreibung: Die Tatsache, dass der Kanal virtuell ist, macht seine
Konstruktion in einer digital organisierten Welt nicht unwirklich.

\section{Drei verschiedene Zeiten}

\begin{definition}[Hostordnung]
Die \emph{Hostordnung} \(H=(E_H,<_{H})\) ist die Erzeugungsordnung aller
tatsächlich ausgeführten Ereignisse. Für den vorliegenden Zeugen ist
\(E_H=\{h_1,\ldots,h_9\}\) und \(h_i<_{H}h_j\) genau dann, wenn \(i<j\).
Neue Hostereignisse dürfen nur am Ende angefügt werden.
\end{definition}

\begin{definition}[Virtuelle Zeitadresse]
Eine \emph{virtuelle Zeitadresse} \(v\in V\) bezeichnet einen
reproduzierbaren Zustandsbezug innerhalb eines Modells. Sie kann kleiner als
die virtuelle Gegenwartsadresse sein, obwohl der adressierende Hostevent
später erzeugt wird. Die Adresse ist eine Koordinate des Modells, keine
Behauptung über den physikalischen Erzeugungszeitpunkt.
\end{definition}

\begin{definition}[Effektordnung]
Die \emph{Effektordnung} \(E\) ordnet die geprüften Freigaben externer
Wirkungen. Empfang, Berechnung und Freigabe sind verschiedene Ereignisse.
Insbesondere gilt
\[
 \texttt{TRANSPORT\_ACK}\neq\texttt{EFFECT\_ACK}.
\]
Nur ein valider, frischer und vollständig gebundener \(\DONE\)-Record kann
unter vollständiger Mediation eine gewöhnliche Wirkung freigeben.
\end{definition}

Die drei Ordnungen dürfen Relationen besitzen, aber nicht identifiziert
werden. Ein Hostevent bei \(h=4\) darf die virtuelle Adresse \(v=15\)
adressieren, obwohl der virtuelle Gegenwartsindex \(v=30\) lautet. Daraus
folgt lediglich
\[
 h_{\mathrm{create}}=4,
 \qquad
 v_{\mathrm{target}}=15<30=v_{\mathrm{now}}.
\]
Es folgt nicht \(h_{\mathrm{effect}}<h_{\mathrm{create}}\). Im Zeugen gilt
im Gegenteil \(4<5<6<7<8<9\).

\section{Modell des bidirektionalen virtuellen Kanals}

\subsection{Unveränderliche Quelle und Branch}

Sei \(S=(e_1,\ldots,e_n)\) eine endliche, append-only Quellhistorie. Ihre
kanonische Identität sei \(d_S\). Eine vergangene Adresse \(v_t\) wählt einen
reproduzierbaren Snapshot \(S_{v_t}\) oder eine wohldefinierte Projektion
darauf. Ein Replay schreibt nicht in \(S\), sondern erzeugt einen neuen
Branch
\[
 B=\operatorname{Replay}(S_{v_t},p,m),
\]
wobei \(p\) die gebundene Programm- und Policyversion und \(m\) die
Nachricht ist. Der Branch trägt mindestens die Identitäten von Quelle,
Zieladresse, Nachricht, Programm und Policy.

Diese Festlegung unterscheidet QIK--VRT von einer nachträglichen Umschreibung
der dokumentierten Vergangenheit. Die Aussage lautet nicht: ``Die spätere
Nachricht war schon damals vorhanden.'' Sie lautet: ``Die spätere Nachricht
wurde jetzt an einen eindeutig bezeichneten damaligen Modellzustand gerichtet
und erzeugte jetzt einen getrennten Branch.''

\subsection{Session und zwei Richtungen}

Eine Session \(\sigma\) enthält mindestens
\[
 (\mathit{sessionID},\mathit{branchID},v_n,v_t,
  d_S,d_p,d_{\mathrm{policy}}).
\]
Die Anfrage ist eine endliche Bitfolge \(m\in\bytes\) mit Richtung
\(v_n\rightarrow v_t\). Eine terminierende Antwortfunktion \(f\) erzeugt
\(a=f(S_{v_t},p,m)\in\bytes\). Die Antwort trägt dieselbe Sessionbindung und
die Richtung \(v_t\rightarrow v_n'\). Anfrage und Antwort benötigen getrennte
Nachrichtenidentitäten, Sequenznummern, Integritätsnachweise und Effect-Ack-
Entscheidungen.

\subsection{Segmentierung}

Für eine Blockgröße \(c\geq1\) wird \(m\) in
\[
 k=\left\lceil\frac{|m|}{c}\right\rceil
\]
Blöcke zerlegt. Ein vollständiges Manifest bindet Gesamtlänge, Blockanzahl,
Blocknummern, Inhaltsidentitäten, Session, Zieladresse und Gesamtidentität.
Die Reassemblierung akzeptiert nur genau die erwartete Folge. Fehlende,
veränderte, fremde oder widersprüchliche Blöcke verhindern die Freigabe.
Mehrfachzustellung darf höchstens idempotent wirken.

\section{Der ausgeführte ISO-C90-Zeuge}

Der Zeuge liegt als
\texttt{tools/qikvrt\_bidirectional\_virtual\_channel\_witness.c} vor. Er
benötigt außer einer ISO-C90-Standardbibliothek keine Abhängigkeiten. Der
Buildvertrag lautet:

\begin{Verbatim}[fontsize=\small]
cc -std=c90 -pedantic-errors -Wall -Wextra -Werror -O2 \
  tools/qikvrt_bidirectional_virtual_channel_witness.c \
  -o qikvrt_bidirectional_virtual_channel_witness
./qikvrt_bidirectional_virtual_channel_witness
\end{Verbatim}

Die Detailausführung erzeugt folgende Hostspur:

\begin{longtable}{@{}rll@{}}
\toprule
\textbf{Host} & \textbf{Operation} & \textbf{Virtueller Bezug} \\
\midrule
\endhead
\(h=1\) & Quellereignis & \(10\rightarrow10\)\\
\(h=2\) & Quellereignis & \(20\rightarrow20\)\\
\(h=3\) & Quellereignis & \(30\rightarrow30\)\\
\(h=4\) & Anfrage erzeugen & \(30\rightarrow15\)\\
\(h=5\) & Transportquittung Anfrage & \(30\rightarrow15\)\\
\(h=6\) & deterministischer Replay & \(15\rightarrow15\)\\
\(h=7\) & Antwort erzeugen & \(15\rightarrow30\)\\
\(h=8\) & Transportquittung Antwort & \(15\rightarrow30\)\\
\(h=9\) & lokaler Demo-Effect-Commit & \(15\rightarrow30\)\\
\bottomrule
\end{longtable}

Die Anfrage umfasst 257 Byte und 16 Blöcke; die deterministisch daraus
gebildete Antwort umfasst 258 Byte und ebenfalls 16 Blöcke. Die beobachteten
nichtkryptographischen Testdigests lauten:

\begin{center}
\begin{tabular}{@{}lr@{}}
\toprule
Quelle vor / nach & \(41545/41545\)\\
Anfrage gesendet / rekonstruiert & \(53716/53716\)\\
Antwort Replay A / Replay B / empfangen & \(23636/23636/23636\)\\
\bottomrule
\end{tabular}
\end{center}

Die Digestfunktion arbeitet absichtlich nur modulo 65521 und ist kein
kryptographischer Hash. Der eigentliche Identitätstest prüft zusätzlich
Länge und sämtliche Bytes. Die Grenzfallserie verwendet Payloadlängen
0, 1, 16, 17, 18, 31, 255, 256, 257 und 4096. Alle zehn Fälle bestehen in
beiden Richtungen. Ein gezielt ausgelassener zweiter Block wird abgelehnt.

\statusbox{QAmber}{
\textbf{Geltungsgrenze des Laufs.}
Der Zeuge beweist die Existenz und Reproduzierbarkeit dieser endlichen
Ausführung. Er ist weder ein Beweis für unendliche Speicherkapazität noch ein
universeller Beweis über alle C-Implementierungen. Der allgemeine Satz folgt
separat aus der folgenden Konstruktion.
}

\section{Vier konstruktive Sätze}

\begin{satz}[Präfixintegrität]\label{thm:prefix}
Sei \(S\) eine append-only Quellhistorie. Wenn Replay ausschließlich einen
neuen Branch erzeugt und kein ausführbarer Nebenpfad Schreibrechte auf das
Quellpräfix besitzt, dann bleibt jedes vor dem Replay vorhandene Präfix von
\(S\) bytegleich erhalten.
\end{satz}

\begin{proof}
Append-only erlaubt auf \(S\) nur die Operation, neue Elemente hinter dem
bisherigen Ende anzufügen. Branch-only Replay schreibt definitionsgemäß in
ein anderes Objekt \(B\). Nach Voraussetzung existiert kein weiterer
Schreibpfad auf das alte Präfix. Damit gibt es in der Übergangsrelation keine
Operation, die ein Byte dieses Präfixes ändern kann. Induktion über die
endliche Folge der Replay- und Transporteereignisse erhält daher jedes
Präfixbyte. Die Behauptung gilt nicht, wenn ein unmediierter Mutationspfad
existiert; genau dieser Fall ist eine Verletzung der Voraussetzung.
\end{proof}

\begin{satz}[Hostkausalität trotz rückwärtsgerichteter virtueller Adresse]
\label{thm:host}
Sei die Hostsemantik append-only und werde jedes abgeleitete Artefakt erst
nach dem Sendeereignis erzeugt. Dann gilt für jeden Replay-, Antworts- und
Commit-Event \(h_x>h_{\mathrm{send}}\), auch wenn
\(v_{\mathrm{target}}<v_{\mathrm{now}}\).
\end{satz}

\begin{proof}
In einer append-only Hostsemantik erhält jedes neue Ereignis einen Index, der
größer ist als alle bereits vorhandenen Indizes. Replay, Antwort und Commit
werden nach dem Senden erzeugt und sind deshalb neue Ereignisse. Also ist
ihr Hostindex jeweils größer als \(h_{\mathrm{send}}\). Die Relation der
virtuellen Adressen kommt in dieser Schlusskette nicht vor. Sie kann daher
die Hostungleichung weder umkehren noch aufheben.
\end{proof}

\begin{satz}[Jede endliche Bitfolge ist vollständig übertragbar]
\label{thm:finite}
Für jedes \(m\in\bytes\) und jede Blockgröße \(c\geq1\) existiert eine
endliche Segmentierung, die \(m\) unter folgenden Bedingungen bytegenau
rekonstruiert: integre Blöcke, eindeutige Sequenznummern, vollständiges
Manifest, schließlich erfolgreiche Zustellung jedes fehlenden Blocks und
ausreichende endliche Ressourcen für genau diesen Lauf.
\end{satz}

\begin{proof}
Ist \(|m|=0\), ist die leere Folge bereits die eindeutige Rekonstruktion. Für
\(|m|>0\) definiere
\(k=\lceil |m|/c\rceil\) und für \(0\leq i<k\) den Block
\[
 m_i=m[ic:\min((i+1)c,|m|)].
\]
Die Folge ist endlich, weil \(|m|\) und \(c\) endlich und \(c>0\) sind.
Eventual Delivery liefert jeden Block nach endlich vielen Wiederholungen.
Eindeutige Sequenznummern ordnen die Blöcke; Integritätsprüfung verwirft
abweichende Inhalte; das Manifest erkennt fehlende oder fremde Blöcke.
Konkatenation in der Reihenfolge \(0,\ldots,k-1\) ergibt nach Definition
exakt \(m\). Für jeden einzelnen Lauf werden nur \(|m|\) Payloadbytes plus
endliche Metadaten benötigt. Behauptet wird eine parametrisierte Familie für
jede endliche Nachricht, nicht die Speicherung aller möglichen Nachrichten
auf einer einzigen festen Maschine.
\end{proof}

\begin{satz}[Bidirektionale Komposition]\label{thm:bidir}
Seien \(m\in\bytes\), \(v_t<v_n\) und \(f\) eine auf dem adressierten
Zustand und \(m\) terminierende Antwortfunktion mit endlicher Ausgabe. Gelten
die Voraussetzungen aus \cref{thm:finite} für beide Richtungen, sind
Session-, Branch-, Programm- und Policyidentitäten eindeutig gebunden und
sind beide Wirkungspfade vollständig mediiert, dann existiert eine endliche
Ausführung
\[
 m:v_n\longrightarrow v_t,
 \qquad
 f(S_{v_t},m):v_t\longrightarrow v_n',
\]
in der beide Payloads bytegenau rekonstruiert werden und alle erzeugenden
Hostereignisse streng vorwärts geordnet bleiben.
\end{satz}

\begin{proof}
Wende \cref{thm:finite} zuerst auf \(m\) an. Nach endlicher Zeit liegt \(m\)
bytegenau beim Branch \(S_{v_t}\) vor. Die Terminierung von \(f\) liefert
nach endlich vielen Hostschritten eine endliche Antwort \(a\). Wende
\cref{thm:finite} auf \(a\) in Gegenrichtung an. Damit liegt \(a\)
bytegenau an \(v_n'\) vor. Eindeutige Session- und Inhaltsbindungen schließen
eine Vermischung der beiden Richtungen mit einer fremden Session aus.
Vollständige Mediation verhindert ordinary release vor dem jeweils
erforderlichen Gateentscheid. Aus \cref{thm:host} folgt, dass sämtliche
tatsächlichen Erzeugungsereignisse vorwärts geordnet bleiben. Die beiden
gerichteten Kanten schließen daher einen virtuellen Adressdialog, aber keinen
physikalischen Kausalkreis.
\end{proof}

\begin{korollar}[Von Text zu beliebigen endlichen Dateien]
Text, Bild, Audioclip, Programm, Modellgewicht oder wissenschaftlicher
Datensatz unterscheiden sich für \cref{thm:finite} nur durch ihre endliche
Bytefolge und ihre Länge. Die Byteübertragung ist daher nicht auf kurze
Nachrichten beschränkt. Korrektes semantisches Verstehen folgt daraus jedoch
nicht.
\end{korollar}

\section{Byteidentität, Bedeutung und Wahrheit}

Die Formulierung ``vollständige Information'' besitzt mindestens drei
verschiedene Lesarten:

\begin{enumerate}
  \item \textbf{syntaktisch:} Alle erwarteten Bytes sind in der richtigen
        Reihenfolge identisch rekonstruiert;
  \item \textbf{semantisch:} Sender und Empfänger verwenden denselben
        versionsgebundenen Decoder und genügend gemeinsamen Kontext;
  \item \textbf{epistemisch:} Die rekonstruierte Aussage ist wahr oder das
        Modell bildet seinen Gegenstand angemessen ab.
\end{enumerate}

\cref{thm:finite} beweist die erste Lesart. Die zweite benötigt einen
explizit gebundenen Codec, eine Ontologie oder eine ausführbare Semantik. Die
dritte benötigt Evidenz über den Gegenstand. Ein Hash kann Inhalt binden,
aber weder Bedeutung noch Wahrheit erzeugen. Ein vollständig übertragenes
falsches Dokument bleibt falsch; ein bytegleich übertragenes Modell kann
empirisch unzutreffend sein.

\section{EFFECT\_ACK und Wirkungssicherheit}

Transportprotokolle beantworten bewusst engere Fragen als eine
Wirkungsfreigabe. Das End-to-End-Argument zeigt seit langem, dass bestimmte
Korrektheitsfunktionen nur an den Anwendungsenden vollständig geprüft werden
können \cite{saltzer}. TCP liefert einen zuverlässigen geordneten Bytestrom
innerhalb seines Vertrags, aber keine Erlaubnis für eine Zahlung, Publikation
oder Aktuatorbewegung \cite{rfc9293}. HTTP 202 trennt Annahme einer Anfrage
von ihrer abgeschlossenen Verarbeitung \cite{rfc9110}; AMQP und MQTT kennen
ebenfalls differenzierte Zustell- und Verarbeitungszustände
\cite{amqp,mqtt}. \qikvrt{} beansprucht daher nicht die Erfindung der ersten
``semantischen Quittung''.

Der engere Beitrag des Effect-Ack-Vertrags ist eine fail-closed
Autorisierungsentscheidung für eine konkrete Wirkung. Sie bindet Input,
Policy, Evidenz, Provenienz, Verantwortung, Frische, Vorgängerzustand und
beobachtetes Ergebnis. Der bestehende Lean-Kern prüft neun benannte Sätze:

\begin{itemize}
  \item \texttt{release\_eq\_true\_iff};
  \item \texttt{release\_requires\_valid\_past};
  \item \texttt{release\_requires\_valid\_future};
  \item \texttt{release\_requires\_effect\_ack};
  \item \texttt{future\_boundary\_is\_counterfactually\_relevant};
  \item \texttt{future\_boundary\_does\_not\_overwrite\_past};
  \item \texttt{identifier\_bound\_eq\_true\_iff};
  \item \texttt{reciprocal\_closure\_eq\_true\_iff};
  \item \texttt{reciprocal\_closure\_requires\_cause\_and\_effect}.
\end{itemize}

Diese Sätze gelten im endlichen booleschen Strukturmodell des gebundenen
Kernel-Receipts. Insbesondere ist die Nichtüberschreibung der
Vergangenheitsprojektion definitional; sie beweist nicht von selbst die
Schreibschutzkonfiguration eines realen Speichers. Die kontrafaktische
Zukunftsrelevanz vergleicht ein konkretes zugelassenes mit einem konkret
abgelehnten Modellobjekt; sie ist keine Beobachtung aus der physikalischen
Zukunft.

\begin{satz}[DONE-only unter vollständiger Mediation]
Wenn jeder Pfad zu einem Executor durch dasselbe Gate führt und ordinary
release genau dann zugelassen wird, wenn der Record vollständig validiert
ist und den Zustand \(\DONE\) trägt, kann kein Non-DONE-Record ordinary
release auslösen.
\end{satz}

\begin{proof}
Nach Voraussetzung besitzt der Executor nur eine eingehende gewöhnliche
Freigabekante. Deren Prädikat enthält \(\mathit{valid}(r)\land
\mathit{state}(r)=\DONE\). Für jeden Non-DONE-Record ist die zweite Konjunkte
falsch, also die ganze Freigabebedingung falsch. Ein Bypass würde die
Mediationsvoraussetzung falsifizieren und darf nicht als Gegenbeispiel zum
bedingten Satz umetikettiert werden.
\end{proof}

\subsection{Exactly-once ist ein eigener Vertrag}

EFFECT\_ACK allein macht einen nichttransaktionalen externen Aktuator nicht
``exactly once''. At-most-once benötigt einen atomaren Dedup-/Effect-Commit
oder einen idempotenten Executor. Liveness benötigt zusätzlich eventual
delivery und eventual availability. Erst die Verbindung beider Seiten
rechtfertigt unter den benannten Fehlerannahmen exactly-once. Ein Crash nach
realer Wirkung, aber vor persistierter Quittung bleibt ohne atomare Kopplung
ein klassisches Unsicherheitsfenster.

\section{Verwandte Arbeiten und Neuheitsgrenze}

Keine der folgenden Grundideen wird als erstmalig beansprucht:

\begin{longtable}{@{}L{0.28\linewidth}L{0.29\linewidth}L{0.35\linewidth}@{}}
\toprule
\textbf{Vorarbeit} & \textbf{Bereits bekannt} & \textbf{Scope-Unterschied} \\
\midrule
\endhead
Lamport \cite{lamport} & Ereignisordnung und Uhren sind zu trennen. & Drei
explizite Ordnungen plus Wirkungsgate.\\
Jefferson / Time Warp \cite{jefferson} & Kleinere virtuelle Zeitstempel,
Rollback und Antimessages. & Unveränderte Quelle, neuer Branch und
Freigabegate statt Quellrollback.\\
Bitemporale Datenbanken \cite{snodgrass} & Valid time und transaction time. &
Ausführung, Sessionreziprozität und Effektordnung.\\
Retroaktive Datenstrukturen \cite{demaine} & Operationen an vergangenen
logischen Zeiten. & Keine Umschreibung der Basishistorie; gebundener Branch.\\
Event Sourcing \cite{fowler} & Append-only Log und Rekonstruktion früherer
Zustände. & Typisierte Zieladresse und getrennte Wirkungsautorisierung.\\
End-to-End \cite{saltzer} & Transporterfolg ist keine
Anwendungskorrektheit. & Geschlossener Policy-/Evidenz-/Verantwortungsrecord.\\
Effektsysteme \cite{lucassen,plotkin} & Effekte werden typisiert und
interpretiert. & Dynamische Autorisierung einer konkreten externen Wirkung.\\
Proof-Carrying Code / Authorization \cite{necula,appel} & Beweise können
Sicherheit oder Berechtigung tragen. & Temporale Replaybindung plus
antizipierte und beobachtete Wirkung.\\
Reversible Berechnung \cite{bennett,danos} & Logische oder kausal konsistente
Umkehr ist formalisiert. & QIK--VRT macht die Quelle nicht rückgängig, sondern
erzeugt einen neuen Branch.\\
\bottomrule
\end{longtable}

Die Einzelmechanismen besitzen somit erhebliche Vorarbeiten. Der hier
untersuchte Neuheitskandidat ist ihre explizite Verbindung:

\begin{quote}
Vergangenheitsadressiertes deterministisches Branch-Replay über einer
unveränderlichen Quellhistorie, getrennte Host-, virtuelle Ziel- und
Effektordnungen, bidirektionale Sessionkorrelation und eine fail-closed
Wirkungsfreigabe, die Input, Policy, Evidenz, Provenienz, Verantwortung,
Frische und beobachtete Wirkung bindet.
\end{quote}

In der gezielten Literaturprüfung wurde kein hier zitiertes System gefunden,
das diese vollständige Verbindung spezifiziert. Das ist ein scope-gebundener
Forschungsneuheitsanspruch, keine Behauptung weltweiter Priorität. Eine solche
Priorität würde systematische Literatur-, Patent- und Standardrecherche sowie
unabhängiges Peer Review verlangen.

\section{Beispiele}

\subsection{Softwarefehler}

Am Freitag wird eine Sicherheitsregel entdeckt. Der unveränderliche
Montagssnapshot einer Software wird durch seinen Digest adressiert. Die neue
Regel wird als Anfrage an die virtuelle Montagsadresse geschickt. Ein Branch
reproduziert den alten Zustand, wendet die Regel an und antwortet: ``Dieser
Angriff wäre blockiert worden'' oder ``Die Regel reicht nicht aus''. Die
reale Montagsausführung ändert sich nicht. Das heutige Verständnis und eine
heute freizugebende Reparatur können sich sehr wohl ändern.

\subsection{Digitaler Zwilling eines Fahrzeugs}

Ein gespeicherter Zustand enthält Softwareversion, Sensorlage, Wetter und
Entscheidungspfad. Eine später entdeckte Gegenlichtregel wird an diesen
Zustand adressiert. Der Branch antwortet mit der alternativen Trajektorie.
Vor einer realen Aktualisierung prüft EFFECT\_ACK Quelle, Modellversion,
Risiko und Verantwortungsfreigabe. Eine bloße erfolgreiche Simulation darf
keine reale Fahrzeugflotte verändern.

\subsection{Medizinische Reanalyse}

Ein neues Auswertungsmodell kann rechtmäßig gespeicherte frühere Daten
reanalysieren. Die Originalbefunde bleiben unverändert; die neue Analyse
bildet einen getrennten Branch. Bytegenaue Übertragung des Modells und der
Daten garantiert noch keine medizinische Richtigkeit. Dafür bleiben
Validierung, Zulassung, Kontext, Datenschutz und ärztliche Verantwortung
eigene Gates.

\subsection{Eine Anweisung, die trotz Empfang blockiert bleibt}

Eine Maschine empfängt alle Bytes einer Steueranweisung und bestätigt den
Transport. Der Policy-Digest stimmt jedoch nicht mit dem autorisierten
Branch überein. Der Record kann daher nicht DONE sein. Vollständige Mediation
blockiert die reale Wirkung, obwohl die Nachricht technisch angekommen ist.
Dieses Beispiel zeigt, warum bidirektionaler Transport und
bidirektional verantwortete Wirkung verschiedene Leistungen sind.

\section{Warum daraus keine physikalische Zeitmaschine folgt}

Ein Computer ist selbstverständlich ein physikalisches System. Daraus folgt
aber nicht, dass jede interne Koordinate des Computers dieselbe physikalische
Bedeutung besitzt wie ihre Bezeichnung. Die virtuelle Adresse \(v=15\) ist
eine Modelladresse. Damit sie einen früheren Raumzeitpunkt bezeichnete,
müsste eine zusätzliche operative Abbildung
\[
 \Phi:V\longrightarrow\mathcal{E}_{\mathrm{spacetime}}
\]
definiert und experimentell bestätigt werden. Keine Zeile des C90-Zeugen
liefert diese Abbildung.

Ein Gegenmodell reicht aus, um den Fehlschluss sichtbar zu machen: Im Zeugen
gilt die virtuelle Kante \(30\rightarrow15\), während alle Hostereignisse
streng \(1\rightarrow2\rightarrow\cdots\rightarrow9\) laufen. Das Modell
erfüllt virtuelle Retroadressierung und verletzt zugleich keine
Vorwärtskausalität des Hosts. Daher impliziert virtuelle
Retroadressierung allein keine physikalische Rückwärtssignalisierung.

Zeit-symmetrische, all-at-once und nichtklassisch kausal geordnete Modelle
sind seriöse Forschungsgegenstände \cite{wharton,ocb}. Auch
Delayed-Choice-Experimente untersuchen nichtklassische zeitliche
Beschreibungen \cite{ma}. Diese Literatur liefert jedoch keinen
kontrollierbaren QIK--VRT-Kanal, der eine später gewählte Nachricht in einen
bereits versiegelten früheren klassischen Record schreibt. Anschlussfähigkeit
ist keine Identität.

\section{Ein falsifizierbarer Test für die offene Physik}

Eine physikalische Brückenbehauptung muss vor der Messung mehr leisten als
eine Deutung im Nachhinein. Ein minimales Prüfprotokoll benötigt:

\begin{enumerate}
  \item einen früheren Messwert \(Y\) bei \(t_1\), der vor jeder späteren
        Wahl irreversibel und extern prüfbar versiegelt wird;
  \item eine erst bei \(t_2>t_1\) erzeugte, unabhängige randomisierte
        Intervention \(X\);
  \item vorregistrierte Ausschlüsse für Informationslecks, gemeinsame
        Ursachen, Uhrfehler, Postselektion, optionale Stopregeln und
        Mehrfachtests;
  \item eine vorab angegebene quantitative Abweichung und Fehlerkontrolle;
  \item unabhängige Replikation.
\end{enumerate}

Die operative Frage lautet, ob nach allen Kontrollen
\[
 P(Y\mid\operatorname{do}(X=x),\Lambda)
\]
reproduzierbar von \(x\) abhängt, obwohl \(Y\) zuvor versiegelt wurde. In
informationstheoretischer Form müsste die kontrollierbare Kanalkapazität
\[
 C_{X\rightarrow Y}=\max_{P(X)} I(X;Y\mid\Lambda)
\]
positiv und gegen Leckmodelle robust sein \cite{shannon,pearl}. Ein einzelner
p-Wert, ein konditioniertes Teilensemble oder ein später neu berechneter
virtueller Branch genügt nicht. Für physisch bidirektionale Zeitkommunikation
wären zwei entsprechend positive operationale Richtungen erforderlich.

\statusbox{QRed}{
\textbf{Aktueller Befund.}
Für \(C_{X\rightarrow Y}>0\) in dieser physikalischen Bedeutung liefert die
vorliegende Arbeit keine Evidenz. Die offene Brücke ist keine kleine
Programmieraufgabe nach dem virtuellen Beweis, sondern eine logisch
zusätzliche Naturhypothese.
}

\section{Sicherheits- und Fehlermodell}

Ein realer bidirektionaler virtueller Kanal muss mindestens gegen folgende
Fehler fail-closed sein:

\begin{itemize}
  \item \textbf{Branch-Verwechslung:} Richtige Bytes im falschen Snapshot
        sind kein richtiger Dialog.
  \item \textbf{Replay-Angriff:} Ein alter DONE-Record darf keine neue
        Session oder widerrufene Policy autorisieren.
  \item \textbf{Codec-Drift:} Byteidentität bei verschiedener
        Decoder-Version kann verschiedene Bedeutung erzeugen.
  \item \textbf{Quellmutation:} Ein veränderter Basissnapshot zerstört die
        behauptete Gegenweltbindung.
  \item \textbf{Unvollständige Mediation:} Jeder Bypass zum Executor
        widerlegt Deployment-Konformität.
  \item \textbf{Ressourcenerschöpfung:} ``Jede endliche Nachricht'' bedeutet
        nicht, dass jede einzelne Maschine jede Größe akzeptieren muss.
        Quoten und Abbruchregeln sind notwendig.
  \item \textbf{Quittungslücke:} Crash zwischen externer Wirkung und
        Receipt-Persistenz verhindert ohne atomare Kopplung exactly-once.
  \item \textbf{Evidenzillusion:} Ein Digest bindet Bytes, validiert aber
        nicht die Wahrheit der referenzierten Evidenz.
\end{itemize}

\section{Protokollskizze für eine spätere Interoperabilitätsfassung}

Eine standardisierbare Session benötigt mindestens folgende typisierte
Objekte:

\begin{longtable}{@{}L{0.26\linewidth}L{0.66\linewidth}@{}}
\toprule
\textbf{Objekt} & \textbf{Pflichtbindung} \\
\midrule
\endhead
Session & Session-ID, Erzeuger, Frische, Ablaufgrenze, beide virtuelle
Endpunkte.\\
Branch & Quellsnapshot, Zieladresse, Programm-/Codec-/Policy-Digest.\\
Message & Richtung, Sequenz, Gesamtlänge, Blockmanifest, Inhaltsidentität.\\
Replay receipt & Eingaben, deterministische Ausgabe, Laufzeitgrenze,
Provenienz.\\
EFFECT\_ACK request & separate Entscheidung über Verarbeitung am
vergangenen virtuellen Zustand.\\
EFFECT\_ACK response & separate Entscheidung über Wirkung am gegenwärtigen
Zustand.\\
Closure receipt & Korrelation beider Richtungen, beobachtete Wirkung und
verantwortlicher Principal.\\
\bottomrule
\end{longtable}

Das öffentlich vorliegende
\texttt{draft-lohmann-qikvrt-effect-ack-02} spezifiziert den bestehenden
Version-1-Effect-Ack-Record als individuelles Internet-Draft
\cite{effectack}. Es ist kein RFC und keine IETF-Billigung. Das vorliegende
wissenschaftliche Paper ist nicht selbst ein Internet-Draft. Erst ein
separat geprüfter Revisionskandidat dürfte entscheiden, ob die Sessionobjekte
nur ein informatives Anwendungsprofil sind oder eine neue Wire-Version mit
CDDL, Sicherheitsanalyse und Interoperabilitätsvektoren erfordern.

\section{Reproduzierbarkeit und aktuelle Grenzen}

Der neue C90-Zeuge wurde lokal mit strikten Warnungen und
\texttt{-pedantic-errors} gebaut und ausgeführt. XeLaTeX und Poppler stehen
für den dokumentierten PDF-Build zur Verfügung. Der kanonische Lean-4.19-
Runtime war in der aktuellen lokalen Umgebung nicht vorhanden; deshalb wird
kein neuer Kernel-Receipt erfunden. Die früheren neun CTM-Theoreme bleiben
durch ihren vorhandenen Receipt gebunden, die neuen
\cref{thm:prefix,thm:host,thm:finite,thm:bidir} tragen bis zu einem frischen
Kernel-Lauf den Status \texttt{CONDITIONAL/PAPER-PROVED} und in der
maschinenlesbaren Zenodo-Policy fail-closed \texttt{OPEN}.

Diese doppelte Bezeichnung ist beabsichtigt: Ein vollständiger Papierbeweis
ist wissenschaftliche Mathematik; die Repository-Policy reserviert
\texttt{FORMAL\_PROVED} jedoch strenger für kernelgeprüfte Claims. Der
fehlende Kernellauf entwertet den Beweis nicht, darf aber auch nicht
verschwiegen werden.

\section{Bedrohungen der Gültigkeit}

\begin{itemize}
  \item Der C90-Zeuge benutzt einen bewusst kleinen deterministischen
        Antworttransform und keine unabhängige Netzwerkimplementierung.
  \item Der Längensweep ist endlich; Universalität stammt aus dem
        Konstruktionsbeweis, nicht aus dem Sweep.
  \item Die Digestwerte des Zeugen sind Testchecksummen, keine
        Kollisionssicherheitsbehauptung.
  \item Eventual Delivery, ausreichende Ressourcen, deterministisches Replay
        und vollständige Mediation sind starke, explizite Voraussetzungen.
  \item Eine terminierende Antwortfunktion ist nicht für jedes beliebige
        Programm garantiert; der Satz löst das Halteproblem nicht.
  \item Weltweite Neuheit ist nicht durch eine gezielte Literaturauswahl
        bewiesen.
  \item Repository-Gates und maschinelle Receipts ersetzen kein unabhängiges
        Peer Review.
  \item Virtuelle Nützlichkeit und physikalische Retrokausalität sind
        logisch verschiedene Behauptungen.
\end{itemize}

\section{Claim-Matrix in lesbarer Form}

\begin{longtable}{@{}L{0.13\linewidth}L{0.22\linewidth}L{0.57\linewidth}@{}}
\toprule
\textbf{ID} & \textbf{Disposition} & \textbf{Aussage und Grenze} \\
\midrule
\endhead
VTI-001 & EXECUTED & Der neue C90-Lauf schließt Anfrage und Antwort zwischen
\(v=30\) und \(v=15\) bei strikt wachsendem Hostclock; nur dieser endliche
Zeuge.\\
VTI-002 & EXECUTED & Zehn Payloadgrenzen bis 4096 Byte bestehen in beiden
Richtungen; ein fehlender Block wird verworfen; kein Universalbeweis.\\
VTI-003 & CONDITIONAL & Append-only Quelle plus branch-only Replay erhält das
Quellpräfix; mutierende Nebenpfade sind ausgeschlossen.\\
VTI-004 & CONDITIONAL & Monotone Hostsemantik schließt einen vor dem Senden
erzeugten Replayeffekt aus; keine physikalische Rückwirkung.\\
VTI-005 & CONDITIONAL & Jede einzelne endliche Bitfolge ist unter den
angegebenen Liefer-, Integritäts- und Ressourcenbedingungen bytegenau
übertragbar.\\
VTI-006 & CONDITIONAL & Zwei endliche Kanäle plus terminierende Antwort und
Gates ergeben einen bidirektionalen virtuellen Dialog.\\
VTI-007 & \shortstack[l]{KERNEL\\VERIFIED} & Bestehender CTM-Kern bindet DONE-only-Freigabe,
Zukunftsbedingung, Vergangenheitsprojektion und reziproke Identitäten im
endlichen Modell.\\
VTI-008 & CONDITIONAL & Non-DONE kann unter vollständiger Mediation keinen
ordinary effect auslösen; Deployment-Vollständigkeit ist Voraussetzung.\\
VTI-009 & OPEN & Exactly-once benötigt atomare Dedup-/Effect-Kopplung oder
Idempotenz plus Liveness; EFFECT\_ACK allein genügt nicht.\\
VTI-010 & NORMATIVE & Byteidentität, semantische Rekonstruktion und Wahrheit
werden als getrennte Prüfziele behandelt.\\
VTI-011 & OPEN & Die vollständige Mechanismenkonjunktion ist ein
scope-gebundener Neuheitskandidat, keine Weltprioritätsbehauptung.\\
VTI-012 & OPEN\_PHYSICAL & Eine Brücke von virtueller Adresse zu früherem
Raumzeitereignis ist unbelegt.\\
VTI-013 & OPEN\_PHYSICAL & Positive rückwärtsgerichtete physikalische
Kanalkapazität verlangt den beschriebenen vorregistrierten Test.\\
\bottomrule
\end{longtable}

\section{Schluss}

Das Ergebnis ist stark genug, ohne seine Kategorie zu wechseln:

\begin{enumerate}
  \item Ein bidirektionaler virtueller Zeitkanal ist nicht nur denkbar,
        sondern in einem strikten ISO-C90-Zeugen ausgeführt.
  \item Der Schritt von kurzen Demo-Payloads zu jeder einzelnen endlichen
        Bitfolge folgt konstruktiv aus Segmentierung, Manifest,
        Integritätsprüfung und Reassemblierung.
  \item Anfrage und Antwort lassen sich zu einem endlichen virtuellen Dialog
        komponieren, wenn Replay und Antwort terminieren und beide Richtungen
        vollständig mediiert sind.
  \item Die Quellhistorie kann unverändert bleiben; die alternative
        Entwicklung ist ein neuer Branch.
  \item Transporterfolg und verantwortete Wirkung bleiben getrennt.
  \item Alle tatsächlichen Hostereignisse laufen vorwärts. Deshalb beweist
        der virtuelle Kanal keine physikalische Rückwärtssignalisierung.
  \item Die physikalische Erweiterung ist klar formulierbar und
        falsifizierbar, aber aktuell offen.
\end{enumerate}

Die Informatik gewinnt damit einen adressierbaren Dialograum über
reproduzierbare Zustände. Die Mathematik liefert die Bedingungen, unter
denen endliche Information vollständig hin- und zurückgeführt werden kann.
Die Physik erhält eine scharf abgegrenzte Brückenfrage statt einer
Begriffsverwechslung.

\begin{center}
\large
\textbf{Heute fragt gestern. Das virtuelle Gestern antwortet dem Heute.}\\
\textbf{Die wirkliche Geschichte bleibt dabei unverändert.}\\[0.6em]
\textit{Quod erat demonstrandum -- im ausgewiesenen virtuellen Scope.}\\
Ingolf Lohmann
\end{center}

\begin{thebibliography}{99}

\bibitem{shannon}
C. E. Shannon,
``A Mathematical Theory of Communication,''
\emph{Bell System Technical Journal}, 27, 379--423 and 623--656, 1948.
\url{https://doi.org/10.1002/j.1538-7305.1948.tb01338.x}

\bibitem{lamport}
L. Lamport,
``Time, Clocks, and the Ordering of Events in a Distributed System,''
\emph{Communications of the ACM}, 21(7), 558--565, 1978.
\url{https://doi.org/10.1145/359545.359563}

\bibitem{jefferson}
D. R. Jefferson,
``Virtual Time,''
\emph{ACM Transactions on Programming Languages and Systems}, 7(3),
404--425, 1985.
\url{https://doi.org/10.1145/3916.3988}

\bibitem{snodgrass}
R. T. Snodgrass and I. Ahn,
``Temporal Databases,''
\emph{IEEE Computer}, 19(9), 35--42, 1986.
\url{https://doi.org/10.1109/MC.1986.1663327}

\bibitem{demaine}
E. D. Demaine, J. Iacono and S. Langerman,
``Retroactive Data Structures,''
\emph{ACM Transactions on Algorithms}, 3(2), Article 13, 2007.
\url{https://doi.org/10.1145/1240233.1240236}

\bibitem{fowler}
M. Fowler,
``Event Sourcing,'' 2005.
\url{https://martinfowler.com/eaaDev/EventSourcing.html}

\bibitem{saltzer}
J. H. Saltzer, D. P. Reed and D. D. Clark,
``End-to-End Arguments in System Design,''
\emph{ACM Transactions on Computer Systems}, 2(4), 277--288, 1984.
\url{https://doi.org/10.1145/357401.357402}

\bibitem{lucassen}
J. M. Lucassen and D. K. Gifford,
``Polymorphic Effect Systems,'' in \emph{POPL '88}, 47--57.
\url{https://doi.org/10.1145/73560.73564}

\bibitem{plotkin}
G. Plotkin and J. Power,
``Algebraic Operations and Generic Effects,''
\emph{Applied Categorical Structures}, 11, 69--94, 2003.
\url{https://doi.org/10.1023/A:1023064908962}

\bibitem{necula}
G. C. Necula,
``Proof-Carrying Code,'' in \emph{POPL '97}, 106--119.
\url{https://doi.org/10.1145/263699.263712}

\bibitem{appel}
A. W. Appel and E. W. Felten,
``Proof-Carrying Authentication,'' in \emph{CCS '99}, 52--62.
\url{https://doi.org/10.1145/319709.319718}

\bibitem{bennett}
C. H. Bennett,
``Logical Reversibility of Computation,''
\emph{IBM Journal of Research and Development}, 17(6), 525--532, 1973.
\url{https://doi.org/10.1147/rd.176.0525}

\bibitem{danos}
V. Danos and J. Krivine,
``Reversible Communicating Systems,'' in \emph{CONCUR 2004}, 292--307.
\url{https://doi.org/10.1007/978-3-540-28644-8_19}

\bibitem{pearl}
J. Pearl,
``Causal Diagrams for Empirical Research,''
\emph{Biometrika}, 82(4), 669--688, 1995.
\url{https://doi.org/10.1093/biomet/82.4.669}

\bibitem{wharton}
K. B. Wharton and N. Argaman,
``Colloquium: Bell's Theorem and Locally Mediated Reformulations of Quantum
Mechanics,'' \emph{Reviews of Modern Physics}, 92, 021002, 2020.
\url{https://doi.org/10.1103/RevModPhys.92.021002}

\bibitem{ocb}
O. Oreshkov, F. Costa and C. Brukner,
``Quantum Correlations with No Causal Order,''
\emph{Nature Communications}, 3, 1092, 2012.
\url{https://doi.org/10.1038/ncomms2076}

\bibitem{ma}
X.-S. Ma, J. Kofler and A. Zeilinger,
``Delayed-Choice Gedanken Experiments and Their Realizations,''
\emph{Reviews of Modern Physics}, 88, 015005, 2016.
\url{https://doi.org/10.1103/RevModPhys.88.015005}

\bibitem{rfc9293}
W. Eddy (ed.),
\emph{Transmission Control Protocol}, RFC 9293, August 2022.
\url{https://doi.org/10.17487/RFC9293}

\bibitem{rfc9110}
R. Fielding, M. Nottingham and J. Reschke,
\emph{HTTP Semantics}, RFC 9110, June 2022.
\url{https://doi.org/10.17487/RFC9110}

\bibitem{amqp}
OASIS,
\emph{Advanced Message Queuing Protocol (AMQP) Version 1.0}, 2012.
\url{https://docs.oasis-open.org/amqp/core/v1.0/}

\bibitem{mqtt}
OASIS,
\emph{MQTT Version 5.0}, 2019.
\url{https://docs.oasis-open.org/mqtt/mqtt/v5.0/mqtt-v5.0.html}

\bibitem{effectack}
I. Lohmann,
\emph{QIK-VRT Effect Acknowledgement: Separating Receipt from Authorization
for Downstream Effect},
Internet-Draft draft-lohmann-qikvrt-effect-ack-02, work in progress, 2026.
\url{https://datatracker.ietf.org/doc/draft-lohmann-qikvrt-effect-ack/}

\bibitem{ctmzenodo}
I. Lohmann,
\emph{QIK-VRT und das Effect-Acknowledgement-Protokoll: Kanonischer Speicher
zwischen Vergangenheit und Zukunft}, Zenodo, 2026.
\url{https://doi.org/10.5281/zenodo.21711193}

\end{thebibliography}

\end{document}
